missing label fig:seq-add-user missing label fig:seq-profile missing label fig:seq-login missing label fig:class-sprint1 replaced fig:uc-sprint1 % ============================================================
% CHAPITRE 4 — Sprint 1 : Utilisateurs, profils, authentification
% ============================================================
\chapter{Sprint 1 --- Gestion des Utilisateurs et Profils, Authentification}

\section{Introduction}

Ce premier sprint de développement marque le démarrage concret de la plateforme
\textbf{DAAM}. Il se concentre sur les fondations liées à l'identité numérique~:
authentification sécurisée, gestion des profils et administration des comptes
utilisateurs. Ces briques sont indispensables avant d'aborder le métier du
recrutement, car chaque espace (Admin, RH, Responsable de recrutement,
Candidat) repose sur une connexion fiable et sur des droits clairement
définis.

Durant ce sprint, le microservice \textbf{User}, l'intégration JWT ainsi que
les premières interfaces Angular (connexion, inscription, profil, gestion
des utilisateurs) ont été mis en place.


\section{Objectifs du Sprint}

Les principaux objectifs de ce sprint sont~:
\begin{itemize}
  \item mettre en œuvre l'authentification des utilisateurs via
        Spring~Security et JWT~;
  \item permettre l'inscription d'un candidat et la connexion de tous les
        rôles prévus~;
  \item créer les composants Angular nécessaires à la connexion, à
        l'inscription et à l'affichage / modification du profil~;
  \item permettre à l'administrateur de gérer les comptes (création,
        modification, activation / désactivation, suppression)~;
  \item concevoir les interfaces frontend correspondantes et brancher le
        frontend sur l'API Gateway.
\end{itemize}


\section{Spécification et Analyse des Besoins}

Cette section détaille les exigences fonctionnelles et les scénarios
d'utilisation qui ont guidé le développement durant ce sprint. Le Sprint
Backlog ci-dessous décompose les User Stories en tâches techniques, tandis
que le diagramme de cas d'utilisation illustre les interactions clés des
acteurs avec le système.


\subsection{Sprint Backlog}

Le tableau suivant présente le Sprint Backlog du sprint~1, détaillant les
fonctionnalités, les scénarios utilisateurs, les tâches associées ainsi que
leurs estimations en charge de travail.

\begin{table}[htbp]
  \centering
  \scriptsize
  \caption{Sprint Backlog du Sprint 1}
  \label{tab:sprint1-backlog}
  \begin{tabular}{|p{2.1cm}|c|p{4.0cm}|c|p{3.4cm}|c|}
    \hline
    \textbf{Fonctionnalité} & \textbf{ID} & \textbf{Scénario utilisateur} &
    \textbf{ID Tâche} & \textbf{Tâches} & \textbf{Est.} \\
    \hline
    Authentification
    & US-01
    & En tant qu'utilisateur, je m'authentifie pour accéder à mon espace selon mon rôle.
    & T1
    & Backend login / register + JWT
    & 2 j \\
    \cline{4-6}
    & & & T2 & Interfaces connexion / inscription & 1,5 j \\
    \cline{4-6}
    & & & T3 & Intercepteur HTTP + guards Angular & 1 j \\
    \hline
    Profil
    & US-02
    & En tant qu'utilisateur connecté, je consulte et mets à jour mon profil (y compris photo).
    & T4
    & API \texttt{/me} + upload image
    & 1,5 j \\
    \cline{4-6}
    & & & T5 & Interface profil & 1 j \\
    \hline
    Gestion des utilisateurs
    & US-03
    & En tant qu'Admin, je crée, modifie, active/désactive et supprime des comptes (rôles inclus).
    & T6
    & CRUD utilisateurs backend
    & 2 j \\
    \cline{4-6}
    & & & T7 & Interface admin \texttt{/admin/users} & 1,5 j \\
    \hline
  \end{tabular}
\end{table}


\subsection{Diagramme de cas d'utilisation du Sprint}

Le diagramme de cas d'utilisation du sprint~1 met en évidence les acteurs
principaux (Administrateur, RH, Responsable de recrutement, Candidat /
utilisateur) et les cas liés à l'authentification, au profil et à la gestion
des comptes.

\begin{figure}[htbp]
  \centering
  \reportfig{img/uc-sprint1.png}
  \caption{Diagramme de cas d'utilisation du Sprint 1}
  \label{fig:uc-sprint1}
\end{figure}


\section{Description textuelle des cas d'utilisation}

\subsection{Description textuelle du cas d'utilisation <<~S'authentifier~>>}

Ce cas d'utilisation décrit le processus par lequel un utilisateur accède à
la plateforme de manière sécurisée. Il couvre la validation des identifiants
et la gestion des accès via un token JWT.

\begin{table}[htbp]
  \centering
  \small
  \caption{Description textuelle du cas d'utilisation <<~S'authentifier~>>}
  \label{tab:uc-login}
  \begin{tabular}{|p{3.0cm}|p{10.2cm}|}
    \hline
    \textbf{Titre} & S'authentifier \\
    \hline
    \textbf{Acteur} & Utilisateur (Admin, RH, Responsable de recrutement ou Candidat) \\
    \hline
    \textbf{Précondition} & L'utilisateur possède un compte valide et actif. \\
    \hline
    \textbf{Postcondition} & L'utilisateur est authentifié, reçoit un JWT et est redirigé vers son espace. \\
    \hline
    \textbf{Scénario nominal}
    &
    \begin{enumerate}[leftmargin=*,itemsep=0.15em,topsep=0.2em]
      \item L'utilisateur ouvre la page de connexion.
      \item Le système affiche le formulaire (identifiant / e-mail et mot de passe).
      \item L'utilisateur saisit ses identifiants et valide.
      \item La requête est envoyée via l'API Gateway vers le User Service.
      \item Le service vérifie les identifiants en base de données.
      \item Si les identifiants sont valides, un token JWT est généré.
      \item Le token est renvoyé au client et stocké (localStorage).
      \item L'utilisateur est redirigé selon son rôle
            (Admin~$\rightarrow$ utilisateurs, RH~$\rightarrow$ recrutements,
            Candidat~$\rightarrow$ offres).
    \end{enumerate}
    \\
    \hline
    \textbf{Alternatives}
    & Identifiants invalides, compte désactivé ou verrouillé~: message d'erreur, pas de token. \\
    \hline
  \end{tabular}
\end{table}


\subsection{Description textuelle du cas d'utilisation <<~Modifier profil~>>}

\begin{table}[htbp]
  \centering
  \small
  \caption{Description textuelle du cas d'utilisation <<~Modifier profil~>>}
  \label{tab:uc-profile}
  \begin{tabular}{|p{3.0cm}|p{10.2cm}|}
    \hline
    \textbf{Titre} & Modifier profil \\
    \hline
    \textbf{Acteur} & Utilisateur authentifié \\
    \hline
    \textbf{Précondition} & L'utilisateur est connecté (JWT valide). \\
    \hline
    \textbf{Postcondition} & Les informations du profil sont mises à jour. \\
    \hline
    \textbf{Scénario nominal}
    &
    \begin{enumerate}[leftmargin=*,itemsep=0.15em,topsep=0.2em]
      \item L'utilisateur ouvre la page profil de son espace.
      \item Le système charge les données via \texttt{GET /api/auth/me}.
      \item L'utilisateur modifie les champs autorisés (coordonnées, photo, etc.).
      \item Le client envoie \texttt{PUT /api/auth/me} (et éventuellement l'upload d'image).
      \item Le User Service persiste les modifications.
      \item Le système confirme la mise à jour à l'utilisateur.
    \end{enumerate}
    \\
    \hline
  \end{tabular}
\end{table}


\subsection{Description textuelle du cas d'utilisation <<~Ajouter un compte utilisateur~>>}

\begin{table}[htbp]
  \centering
  \small
  \caption{Description textuelle du cas d'utilisation <<~Ajouter un compte utilisateur~>>}
  \label{tab:uc-add-user}
  \begin{tabular}{|p{3.0cm}|p{10.2cm}|}
    \hline
    \textbf{Titre} & Ajouter un compte utilisateur \\
    \hline
    \textbf{Acteur} & Administrateur \\
    \hline
    \textbf{Précondition} & L'administrateur est authentifié avec le rôle \texttt{ROLE\_ADMIN}. \\
    \hline
    \textbf{Postcondition} & Un nouveau compte est créé avec le rôle choisi. \\
    \hline
    \textbf{Scénario nominal}
    &
    \begin{enumerate}[leftmargin=*,itemsep=0.15em,topsep=0.2em]
      \item L'Admin ouvre la page de gestion des utilisateurs.
      \item Il saisit les informations du compte et sélectionne un rôle
            (RH, Responsable de recrutement, Candidat, etc.).
      \item Le frontend envoie \texttt{POST /api/auth/users}.
      \item Le User Service valide les données, encode le mot de passe et enregistre l'utilisateur.
      \item Le nouveau compte apparaît dans la liste.
    \end{enumerate}
    \\
    \hline
    \textbf{Alternatives}
    & E-mail / username déjà existant~: création refusée avec message d'erreur. \\
    \hline
  \end{tabular}
\end{table}


\subsection{Gestion des autres actions liées aux comptes utilisateurs}

En complément des cas détaillés ci-dessus, le sprint couvre également~:
\begin{itemize}
  \item \textbf{Inscription publique}~: un candidat crée son propre compte
        (\texttt{POST /api/auth/register})~; le rôle attribué est
        \texttt{ROLE\_USER}.
  \item \textbf{Lister les utilisateurs}~: l'Admin consulte l'ensemble des
        comptes (\texttt{GET /api/auth/users}).
  \item \textbf{Modifier un compte}~: l'Admin met à jour les informations
        d'un utilisateur (\texttt{PUT /api/auth/users/\{userId\}}).
  \item \textbf{Activer / désactiver}~: l'Admin change le statut d'un compte
        (\texttt{PATCH /api/auth/users/\{userId\}/status}).
  \item \textbf{Supprimer un compte}~: l'Admin retire un utilisateur
        (\texttt{DELETE /api/auth/users/\{userId\}}).
\end{itemize}


\section{Conception}

\subsection{Diagramme de classes}

Le diagramme de classes du sprint~1 concentre le domaine d'identité autour
de l'entité \textbf{User} (identifiants, coordonnées, rôle, statut, photo de
profil) et des services associés (authentification JWT, gestion des comptes,
stockage de fichiers).

\begin{figure}[htbp]
  \centering
  \reportfig{img/class-sprint1.png}
  \caption{Diagramme de classes du Sprint 1}
  \label{fig:class-sprint1}
\end{figure}


\subsection{Diagrammes de séquence}

\subsubsection{Diagramme de séquence <<~S'authentifier~>>}

\begin{figure}[htbp]
  \centering
  \reportfig{img/seq-login.png}
  \caption{Diagramme de s\'equence --- S'authentifier}
  \label{fig:seq-login}
\end{figure}

\subsubsection{Diagramme de séquence <<~Modifier profil~>>}

\begin{figure}[htbp]
  \centering
  \reportfig{img/seq-profile.png}
  \caption{Diagramme de s\'equence --- Modifier profil}
  \label{fig:seq-profile}
\end{figure}

\subsubsection{Diagramme de séquence <<~Ajouter un compte utilisateur~>>}

\begin{figure}[htbp]
  \centering
  \reportfig{img/seq-add-user.png}
  \caption{Diagramme de s\'equence --- Ajouter un compte utilisateur}
  \label{fig:seq-add-user}
\end{figure}


\section{Réalisation}

Cette section présente les interfaces réalisées pour le sprint~1.

\subsection{Interface d'authentification (login / inscription)}

L'écran d'authentification permet la connexion et l'inscription. Après succès,
le JWT est conservé côté client et la navigation est orientée selon le rôle.

\begin{figure}[htbp]
  \centering
  \reportfig{img/ui-login.png}
  \caption{Interface de connexion / inscription de la plateforme DAAM}
  \label{fig:ui-login}
\end{figure}


\subsection{Interface de gestion des utilisateurs}

L'espace Admin propose la liste des comptes, la création d'utilisateurs avec
choix du rôle, ainsi que les actions de modification, d'activation et de
suppression.

\begin{figure}[htbp]
  \centering
  \reportfig{img/ui-users.png}
  \caption{Interface de gestion des utilisateurs (Admin)}
  \label{fig:ui-users}
\end{figure}


\subsection{Interface du profil utilisateur}

Chaque acteur dispose d'une page profil pour consulter et mettre à jour ses
informations personnelles (y compris la photo de profil).

\begin{figure}[htbp]
  \centering
  \reportfig{img/ui-profile.png}
  \caption{Interface du profil utilisateur}
  \label{fig:ui-profile}
\end{figure}


\section{Conclusion}

Le Sprint~1 a permis de poser le socle d'identité de la plateforme DAAM~:
authentification JWT, inscription candidat, gestion des profils et
administration des comptes multi-rôles. Ces fondations sécurisent l'accès aux
espaces et préparent les sprints suivants dédiés à l'organisation
administrative puis au cœur métier du recrutement.

Le chapitre suivant présente le Sprint~2, consacré à l'organisation Admin
(zones, entreprises / agences et affectations RH).

\clearpage
